\documentclass[11pt]{article}

\usepackage[margin=1in]{geometry}
\usepackage{amsmath,amssymb,amsthm}
\usepackage{graphicx}
\usepackage[hidelinks]{hyperref}

\newtheorem{theorem}{Theorem}
\newtheorem{definition}{Definition}
\newtheorem{remark}{Remark}

\newcommand{\N}{\mathbb N}
\newcommand{\eps}{\varepsilon}

\title{No Unbounded Convex \texorpdfstring{$W$}{W}-Hyperbolic Set Has the UAFPP}
\author{Automated Math Discovery Smoke Test}
\date{July 18, 2026}

\begin{document}
\maketitle

\section*{Source and Open Problem}

This packet concerns the open problem on page 44 of:

\begin{quote}\raggedright
Laurentiu Leustean, \emph{Proof mining in metric fixed point theory and
ergodic theory}, arXiv:0903.1529, Oberwolfach Preprints OWP 2009-05.
\end{quote}

The same UAFPP problem is also the final open problem of the original
Kohlenbach--Leustean paper on approximate fixed points in product spaces.
The source asks whether any unbounded convex subset of a $W$-hyperbolic space
can have the uniform approximate fixed point property for all nonexpansive
self-maps.

\begin{center}
\includegraphics[width=\textwidth]{figures/open_problem_crop.png}
\end{center}

\section*{Definitions}

We use the standard notation from the source. A $W$-hyperbolic space is a
metric space $(X,d)$ equipped with a convexity map
\[
  W:X\times X\times[0,1]\to X.
\]
The only axiom needed below is the arc-length axiom
\[
  d\bigl(W(x,y,\lambda),W(x,y,\mu)\bigr)
  =
  |\lambda-\mu|\,d(x,y)
  \qquad(\lambda,\mu\in[0,1]).
\]
A subset $C\subseteq X$ is convex if
$W(x,y,\lambda)\in C$ for all $x,y\in C$ and $\lambda\in[0,1]$.

\begin{definition}[UAFPP]
Let $C$ be a metric space and let $\mathcal F$ be a class of maps
$T:C\to C$. We say that $C$ has the uniform approximate fixed point property
for $\mathcal F$ if for every $\eps>0$ and every $b>0$ there is $D>0$ such
that, for every $x\in C$ and every $T\in\mathcal F$,
\[
  d(x,Tx)\le b
  \quad\Longrightarrow\quad
  \exists x^*\in C\ 
  \bigl(d(x,x^*)\le D\ \text{and}\ d(x^*,Tx^*)<\eps\bigr).
\]
\end{definition}

\section*{Proof Intuition}

The source already observes that a stronger, basepoint-free uniform property
would force boundedness by testing constant maps. UAFPP avoids that simple
obstruction because it only asks for nearby approximate fixed points when the
chosen base point is moved by at most $b$.

The missing obstruction is still one-dimensional. From a base point $a$ choose
a very long geodesic segment inside the unbounded convex set. Collapse every
point of $C$ to this segment according to its distance from $a$, then move one
unit farther along the segment, stopping only at the far endpoint. This radial
map is nonexpansive because both the distance-to-$a$ function and the
arc-length parametrization of the segment are $1$-Lipschitz. Near $a$, it
increases the distance from $a$ by exactly one, so the reverse triangle
inequality forces displacement at least one throughout any prescribed bounded
neighborhood of $a$.

\section*{Main Result}

\begin{theorem}
Let $C$ be a nonempty unbounded convex subset of a $W$-hyperbolic space
$(X,d,W)$. Then $C$ does not have the UAFPP for the class of all nonexpansive
self-mappings $T:C\to C$.

Consequently, the open problem in arXiv:0903.1529 has a negative answer:
there are no unbounded convex subsets of $W$-hyperbolic spaces with the UAFPP
for all nonexpansive self-maps.
\end{theorem}

\begin{proof}
Fix $a\in C$. Suppose, toward a contradiction, that $C$ has the UAFPP for all
nonexpansive self-maps. Apply the definition with $\eps=1$ and $b=1$, and let
$D>0$ be the corresponding radius.

Since $C$ is unbounded, choose $y\in C$ such that
\[
  L:=d(a,y)>D+1.
\]
For $t\in[0,L]$, define the geodesic segment
\[
  \gamma(t):=W(a,y,t/L).
\]
By convexity of $C$, $\gamma([0,L])\subseteq C$, and by the arc-length axiom,
\[
  d(\gamma(s),\gamma(t))=|s-t|
  \qquad (s,t\in[0,L]).
\]

Define $h:[0,\infty)\to[0,L]$ by
\[
  h(r):=\min\{r+1,L\},
\]
and define $T:C\to C$ by
\[
  Tz:=\gamma\bigl(h(d(a,z))\bigr).
\]
The function $h$ is $1$-Lipschitz, and $z\mapsto d(a,z)$ is $1$-Lipschitz.
Therefore, for all $z,z'\in C$,
\[
\begin{aligned}
  d(Tz,Tz')
  &=
  \bigl|h(d(a,z))-h(d(a,z'))\bigr| \\
  &\le |d(a,z)-d(a,z')|
   \le d(z,z').
\end{aligned}
\]
Thus $T$ is nonexpansive. Moreover,
\[
  d(a,Ta)=d(\gamma(0),\gamma(1))=1,
\]
so the UAFPP hypothesis applies at the base point $a$.

We now show that no point in the $D$-neighborhood of $a$ is a $1$-fixed point
of $T$. Let $z\in C$ satisfy $d(a,z)\le D$. Since $D+1<L$, we have
$h(d(a,z))=d(a,z)+1$. Hence
\[
  d(a,Tz)=d(a,z)+1.
\]
By the reverse triangle inequality,
\[
  d(z,Tz)\ge |d(a,Tz)-d(a,z)|=1.
\]
Thus $d(z,Tz)<1$ fails for every $z\in C$ with $d(a,z)\le D$.

This contradicts the UAFPP conclusion for $\eps=1$, $b=1$, the point $a$, and
the nonexpansive map $T$. Therefore no such unbounded convex set $C$ has the
UAFPP.
\end{proof}

\section*{Verification and Novelty Notes}

The proof is purely metric and uses only the arc-length axiom of
$W$-hyperbolic spaces, convexity of $C$, and the reverse triangle inequality.
No computation is used in the proof.

Before promotion, the run indexes were searched for `0903.1529`, proof mining,
metric fixed point theory, UAFPP, approximate fixed point property, and related
fixed-point keywords. No exact prior packet for this source problem was found.
Nearby registry hits concerned other fixed-point or metastability targets.

A bounded web search using the phrases `uniform approximate fixed point
property`, `unbounded convex subsets`, `W-hyperbolic`, `UAFPP`, and the exact
open-problem wording found the source arXiv survey and the original
Kohlenbach--Leustean article, both presenting the problem as open, but did not
find a later solution.

\section*{Human Review Recommendation}

Review as a likely-valid full solution. The main check is whether the
arc-length property quoted above is indeed part of the source's definition of
$W$-hyperbolic space; it is the standard axiom used in this literature. Once
that is confirmed, the nonexpansiveness and contradiction are immediate.

\begin{thebibliography}{9}

\bibitem{Leustean2009}
L. Leustean,
\emph{Proof mining in metric fixed point theory and ergodic theory},
arXiv:0903.1529, Oberwolfach Preprints OWP 2009-05, 2009.

\bibitem{KohlenbachLeustean2007}
U. Kohlenbach and L. Leustean,
\emph{The approximate fixed point property in product spaces},
Nonlinear Analysis: Theory, Methods \& Applications 66 (2007), 806--818.

\end{thebibliography}

\end{document}
